% GENERATED FILE. Edit canonical lessons, then run npm run generate:books.
% canonical-source-hash: fnv1a64:43656288d13f02de
% canonical-lessons: RU-C472-razdrazhat, RU-C472-razdrazheniye, RU-C472-zlitsya, RU-C472-yarost, RU-C472-vspylchivyy

\chapter{Irritation and Fury}
\label{ch:ru-irritation-and-fury}
\hlchaptermodality{472}

\begin{chapteropening}
\textbf{By the end of this chapter:} \emph{I can talk about irritation and anger.}

Complete the last lesson of chapter 472: I can talk about irritation and anger.
\end{chapteropening}

\section[\texorpdfstring{\ru{раздражать} (razdrazhát)}{раздражать (razdrazhát)}]{\ru{раздражать} (razdrazhát) — to irritate}

\label{lesson:RU-C472-razdrazhat}



\begin{quote}
Before the new one: say the Russian for touchy, then the Russian for to forgive.
\end{quote}

\subsection*{You'll want to know: \ru{раздражать}}
\textbf{\ru{раздражать}} (\emph{razdrazhát}) — \textquotedblleft{}to irritate\textquotedblright{}.

\begin{quote}
\ru{Меня} \ru{раздражает} \ru{шум} \ru{за} \ru{стеной}.
\end{quote}

\begin{quote}
\emph{Menyá razdrazháyet shum za stenóy.}
\end{quote}

\begin{quote}
The noise behind the wall irritates me.
\end{quote}

Imperfective; perfective \ru{раздражить}. The person who is irritated goes in the accusative: \emph{\ru{меня} \ru{раздражает}}.

\begin{cousinweb}
\emph{\ru{раздражать}} is built on \emph{\ru{раз}-} and a stem usually linked to the one in \emph{\ru{дразнить}} 'to tease'; EtymDB lists the verb, glossed 'to irritate', without a source.
\end{cousinweb}

\begin{grammarlens}[title={What you've built}]
One more word for telling, explaining and describing at B1.
\end{grammarlens}

\subsection*{Your turn}
\begin{itemize}
\raggedright
  \item \emph{Say it:} \emph{razdrazhát}
  \item \emph{Say it:} \emph{razdrazhát}, once more
  \item \emph{Say it:} a sentence of your own with \emph{razdrazhát}
  \item \emph{Recall:} say the Russian for touchy, then the Russian for to forgive, then say \emph{razdrazhát} again
\end{itemize}

\begin{tcolorbox}[breakable,colback=teal!4,colframe=teal!35!black,title={Before you move on}]
What does \ru{раздражать} mean? (\textquotedblleft{}To irritate\textquotedblright{}.) Say it once more.
\end{tcolorbox}
\section[\texorpdfstring{\ru{раздражение} (razdrazhéniye)}{раздражение (razdrazhéniye)}]{\ru{раздражение} (razdrazhéniye) — irritation}

\label{lesson:RU-C472-razdrazheniye}



\begin{quote}
Before the new one: say the Russian for to forgive, then the Russian for to irritate.
\end{quote}

\subsection*{You'll want to know: \ru{раздражение}}
\textbf{\ru{раздражение}} (\emph{razdrazhéniye}) — \textquotedblleft{}irritation\textquotedblright{}.

\begin{quote}
\ru{Он} \ru{ответил} \ru{с} \ru{раздражением}.
\end{quote}

\begin{quote}
\emph{On otvétil s razdrazhéniyem.}
\end{quote}

\begin{quote}
He answered irritably.
\end{quote}

Neuter. \emph{\ru{С} \ru{раздражением}} — irritably.

\begin{cousinweb}
\emph{\ru{раздражение}} is built on \emph{\ru{раздражать}} 'to irritate' with \emph{-\ru{ение}}; the verb's stem is usually linked to \emph{\ru{дразнить}} 'to tease', and EtymDB has no entry for the noun.
\end{cousinweb}

\begin{grammarlens}[title={What you've built}]
One more word for telling, explaining and describing at B1.
\end{grammarlens}

\subsection*{Your turn}
\begin{itemize}
\raggedright
  \item \emph{Say it:} \emph{razdrazhéniye}
  \item \emph{Say it:} \emph{razdrazhéniye}, once more
  \item \emph{Say it:} a sentence of your own with \emph{razdrazhéniye}
  \item \emph{Recall:} say the Russian for to forgive, then the Russian for to irritate, then say \emph{razdrazhéniye} again
\end{itemize}

\begin{tcolorbox}[breakable,colback=teal!4,colframe=teal!35!black,title={Before you move on}]
What does \ru{раздражение} mean? (\textquotedblleft{}Irritation\textquotedblright{}.) Say it once more.
\end{tcolorbox}
\section[\texorpdfstring{\ru{злиться} (zlítsya)}{злиться (zlítsya)}]{\ru{злиться} (zlítsya) — to be angry}

\label{lesson:RU-C472-zlitsya}



\begin{quote}
Before the new one: say the Russian for to irritate, then the Russian for irritation.
\end{quote}

\subsection*{You'll want to know: \ru{злиться}}
\textbf{\ru{злиться}} (\emph{zlítsya}) — \textquotedblleft{}to be angry\textquotedblright{}.

\begin{quote}
\ru{Папа} \ru{злится}, \ru{когда} \ru{мы} \ru{опаздываем}.
\end{quote}

\begin{quote}
\emph{Pápa zlítsya, kogdá my opázdyvayem.}
\end{quote}

\begin{quote}
Dad gets angry when we're late.
\end{quote}

Imperfective; perfective \ru{разозлиться}. \emph{\ru{Злиться} \ru{на} \ru{кого}} — to be angry with someone. A little more colloquial than \emph{\ru{сердиться}}.

\begin{cousinweb}
\emph{\ru{злиться}} is built on \emph{\ru{злой}} 'angry, wicked', which is inherited from Proto-Slavic \emph{z\ru{ъ}l\ru{ъ}} 'bad, evil' (EtymDB).
\end{cousinweb}

\begin{grammarlens}[title={What you've built}]
One more word for telling, explaining and describing at B1.
\end{grammarlens}

\subsection*{Your turn}
\begin{itemize}
\raggedright
  \item \emph{Say it:} \emph{zlítsya}
  \item \emph{Say it:} \emph{zlítsya}, once more
  \item \emph{Say it:} a sentence of your own with \emph{zlítsya}
  \item \emph{Recall:} say the Russian for to irritate, then the Russian for irritation, then say \emph{zlítsya} again
\end{itemize}

\begin{tcolorbox}[breakable,colback=teal!4,colframe=teal!35!black,title={Before you move on}]
What does \ru{злиться} mean? (\textquotedblleft{}To be angry\textquotedblright{}.) Say it once more.
\end{tcolorbox}
\section[\texorpdfstring{\ru{ярость} (yárost)}{ярость (yárost)}]{\ru{ярость} (yárost) — fury}

\label{lesson:RU-C472-yarost}



\begin{quote}
Before the new one: say the Russian for irritation, then the Russian for to be angry.
\end{quote}

\subsection*{You'll want to know: \ru{ярость}}
\textbf{\ru{ярость}} (\emph{yárost}) — \textquotedblleft{}fury\textquotedblright{}.

\begin{quote}
\ru{В} \ru{ярости} \ru{он} \ru{бросил} \ru{телефон} \ru{на} \ru{пол}.
\end{quote}

\begin{quote}
\emph{V yárosti on brósil telefón ná pol.}
\end{quote}

\begin{quote}
In a rage he threw his phone on the floor.
\end{quote}

Feminine. The adjective is \emph{\ru{яростный}} — furious.

\begin{cousinweb}
\emph{\ru{ярость}} is inherited from Proto-Slavic \emph{jarost\ru{ь}} 'fury', derived from \emph{jar\ru{ъ}} 'furious' (EtymDB), the root also behind \emph{\ru{яркий}} 'bright'.
\end{cousinweb}

\begin{grammarlens}[title={What you've built}]
One more word for telling, explaining and describing at B1.
\end{grammarlens}

\subsection*{Your turn}
\begin{itemize}
\raggedright
  \item \emph{Say it:} \emph{yárost}
  \item \emph{Say it:} \emph{yárost}, once more
  \item \emph{Say it:} a sentence of your own with \emph{yárost}
  \item \emph{Recall:} say the Russian for irritation, then the Russian for to be angry, then say \emph{yárost} again
\end{itemize}

\begin{tcolorbox}[breakable,colback=teal!4,colframe=teal!35!black,title={Before you move on}]
What does \ru{ярость} mean? (\textquotedblleft{}Fury\textquotedblright{}.) Say it once more.
\end{tcolorbox}
\section[\texorpdfstring{\ru{вспыльчивый} (vspýlchivyy)}{вспыльчивый (vspýlchivyy)}]{\ru{вспыльчивый} (vspýlchivyy) — quick-tempered}

\label{lesson:RU-C472-vspylchivyy}



\begin{quote}
Before the new one: say the Russian for to be angry, then the Russian for fury.
\end{quote}

\subsection*{You'll want to know: \ru{вспыльчивый}}
\textbf{\ru{вспыльчивый}} (\emph{vspýlchivyy}) — \textquotedblleft{}quick-tempered\textquotedblright{}.

\begin{quote}
\ru{Мой} \ru{дядя} \ru{добрый}, \ru{но} \ru{вспыльчивый}.
\end{quote}

\begin{quote}
\emph{Moy dyádya dóbryy, no vspýlchivyy.}
\end{quote}

\begin{quote}
My uncle is kind but quick-tempered.
\end{quote}

Built on \emph{\ru{вспылить}} — to flare up.

\begin{cousinweb}
\emph{\ru{вспыльчивый}} is built on \emph{\ru{вспылить}} 'to flare up', from \emph{\ru{пыл}} 'heat, ardour'; EtymDB has no entry, and \emph{\ru{пыл}} is usually traced to the root of \emph{\ru{пылать}} 'to blaze'.
\end{cousinweb}

\begin{grammarlens}[title={What you've built}]
One more word for telling, explaining and describing at B1.
\end{grammarlens}

\subsection*{Your turn}
\begin{itemize}
\raggedright
  \item \emph{Say it:} \emph{vspýlchivyy}
  \item \emph{Say it:} \emph{vspýlchivyy}, once more
  \item \emph{Say it:} a sentence of your own with \emph{vspýlchivyy}
  \item \emph{Recall:} say the Russian for to be angry, then the Russian for fury, then say \emph{vspýlchivyy} again
\end{itemize}

\begin{tcolorbox}[breakable,colback=teal!4,colframe=teal!35!black,title={Before you move on}]
What does \ru{вспыльчивый} mean? (\textquotedblleft{}Quick-tempered\textquotedblright{}.) Say it once more.
\end{tcolorbox}
